\documentclass{article}
\usepackage[utf8]{inputenc}
\usepackage[T1]{fontenc}
\usepackage{geometry}
\geometry{
    a4paper,
    left=2.5cm,
    right=2.5cm,
    top=2.5cm,
    bottom=2.5cm
}
\usepackage{graphicx}
\usepackage{float}
\usepackage{hyperref}
\usepackage[english]{babel}
\usepackage{amsmath}
\usepackage{amssymb}
\usepackage{array}
\usepackage{booktabs}
\usepackage{caption}
\usepackage{subcaption}
\usepackage{enumitem}
\usepackage{fancyhdr}
\usepackage{titlesec}
\usepackage{url}
\usepackage{textcomp}
\usepackage{xcolor}
\usepackage{listings}
\usepackage{multirow}
\usepackage{longtable}
\usepackage{tabularx}
\usepackage{times}

% Set up page style
\pagestyle{fancy}
\fancyhf{}
\rfoot{\thepage}
\renewcommand{\headrulewidth}{0pt}

% Title formatting
\titleformat{\chapter}[display]{\normalfont\bfseries}{\chaptertitle\thechapter}{0pt}{\Huge}
\titlespacing*{\chapter}{0pt}{50pt}{40pt}

% Document info
\title{Template Document}
\author{Academic Support Agent}
\date{\today}

\begin{document}
\maketitle

\maketitle

\begin{resumen}
\lipsum[1-4]
\end{resumen}

% opcional: incluir para tesis en ingl s;
%  en este caso hay que tener el resumen y abstract
%   en ambos idiomas
%\begin{abstract}
%\lipsum[1-4]
%\end{abstract}

\begin{dedicatoria}
Una dedicatoria corta.
\end{dedicatoria}

\begin{thanks}
\lipsum[1-2]
\end{thanks}
{
\hypersetup{linkcolor=black}



}

\pagestyle{fancy}
\fancyhead{} % clear all header fields
\fancyhead[C]{\nouppercase{\rightmark\hfill\leftmark}}
\fancyhead[C]{\nouppercase{\leftmark\hfill\rightmark}}
\fancyfoot{} % clear all footer fields
\fancyfoot[C]{\thepage}



\subfile{cap1}
\subfile{cap2}
\subfile{cap3}
\subfile{conclu}

%Inlcuir cuantos cap tulos estime necesario
% % Content from cap1.tex would be included here

% % Content from cap2.tex would be included here

% % Content from cap3.tex would be included here

% % Content from conclu.tex would be included here


% ver https://www.overleaf.com/learn/latex/Glossaries
% % Content from glosario.tex would be included here
 % opcional

%\nocite{*} %para incluir todos los items de la bibliografia (aun cuando no sean citados) 



% opcional ...
\begin{appendices}
\subfile{anexoA}
\end{appendices}
\end{document}